\section{Introduction}
\label{sec:v27-introduction}
\label{sec:v26-introduction}
\label{sec:v25-introduction}

The geometry of a dispersing billiard can be observed through trajectories,
lengths, collision positions, or probability laws of selected collision
words.  These data have different inverse content.  An exact planar
collision position lies on the unknown boundary and supplies a geometric
sample.  Its signed transverse projection does not supply the missing
normal coordinate.  We study the geometry encoded by the distribution of
the latter observations for long alternating bridges near a closest pair
of strictly convex scatterers.

Our main local result determines both labelled contact jets from the gap
and one signed conditional endpoint law of each contact type at a single
positive offset.  The bridge probability decays exponentially with the
number of flights, but its conditional law retains a nonlinear profile
formed by repeated returns to both contacts.  We construct this profile
as a relative limit on a collar independent of the flight number and
invert it on the geometric image of the billiard observation map.
Neither reflection symmetry, equal contacts, nor supplied curvatures are
required.  The inverse combines cancellation of the unknown flux
amplitude with a finite-remainder analysis of the nonlinear half-line
action.  This gives a triangular reconstruction of the labelled contact
jets at every prescribed finite order.  For connected analytic boundaries,
intrinsic channel matching then recovers complete obstacles and the
Euclidean period lattice under the marked incidence hypotheses stated
below.

We also construct globally consistent estimators in an unregistered
physical position-sensor experiment, with a charged calibration pilot and
only long even bridges.  We give a direct graph-sampling proof available
in that same richer experiment.  Thus the distinction between geometric
sampling and the inverse from transverse laws is part of the theorem
statements, rather than an inference from the flight restriction.  The
local information results separately identify the effects of deleting
residual times and waiting counts.

\subsection{Relative boundary laws}

For a selected closest-pair channel let \(g>0\) be its gap, let
\(\kappa_0,\kappa_1>0\) be its labelled contact curvatures, and put
\[
 c_b=1+g\kappa_b,\qquad
 \gamma=\operatorname{arcosh}\sqrt{c_0c_1}.
\]
If \(W_j(u,v)\) is the stationary length of a \(j\)-flight bridge with
signed transverse endpoints, define
\[
 E_j=W_j-jg,\qquad
 b_j=\frac{-W_{j,uv}}{-W_{j,uv}(0,0)}.
\]
There are strictly convex actions \(S_b\) and positive amplitudes \(B_b\),
normalized by \(S_b(0)=S_b'(0)=0\) and \(B_b(0)=1\), such that
\begin{equation}\label{eq:v25-intro-relative}
 \|E_j-S_0\oplus S_p\|_{C^k}
 +\|b_j-B_0\otimes B_p\|_{C^k}\le C_k\tau^j,
 \qquad p=j\bmod2,\quad 0<\tau<1.
\end{equation}
The order \(k\) is fixed but arbitrary, including the mixed geometric and
offset derivatives stated in Theorem~\ref{thm:v8-main-relative}.  The
endpoint and offset collars do not shrink with \(j\).  The local forward
construction requires neither finite horizon nor global reflection
symmetry.  Its weighted inverse, differentiated operator comparison, and
uniform hyperbolicity margins are proved in Part I.

At excess time \(d>0\) the successful endpoint--residual-time limit has
density proportional to
\[
 B_0(u)B_p(v)\mathbf1_{\{0<r<d-S_0(u)-S_p(v)\}}.
\]
For same-type even bridges, integrating out \(r\) gives a density
\(f_{b,d}(u,v)\) proportional to
\((d-S_b(u)-S_b(v))_+B_b(u)B_b(v)\).
The family of supports as \(d\) varies determines the unsymmetrized
action by a threshold slice.  The density itself carries more information:
only one positive \(d\) is needed.  Proposition~\ref{prop:v27-support-gauge}
gives a support-preserving action change in the functional density class,
showing why one support alone does not supply the same invariant.

\subsection{A single-offset inverse and intrinsic periodic determination}

On a small interior square, write
\[
 R_b(u,v)=\frac{f_{b,d}(u,v)f_{b,d}(0,0)}
                  {f_{b,d}(u,0)f_{b,d}(0,v)}.
\]
All unknown amplitude and normalization factors cancel, leaving
\begin{equation}\label{eq:v26-intro-ratio}
 1-R_b(u,v)=t_b(u)t_b(v),\qquad
 t_b(u)=\frac{S_b(u)}{d-S_b(u)}.
\end{equation}
At a fixed nonzero anchor \(a\), strict convexity gives
\(t_b(a)>0\).  Consequently
\[
 t_b(u)=\frac{1-R_b(u,a)}{\sqrt{1-R_b(a,a)}},\qquad
 S_b(u)=\frac{dt_b(u)}{1+t_b(u)}.
\]
These are identities of continuous densities determined by the law; point
values of a density are not assumed to be individual observations.  The
fixed-anchor formula has a locally Lipschitz inverse in every fixed
interior \(C^M\) norm, uniformly on compact nondegenerate classes.  It
avoids taking derivatives of a square root at the minimum of \(S_b\).

The gap is a separate datum, obtained from physical onset \(jg\).  If
\(a_b=S_b''(0)\), then
\[
 \gamma=\operatorname{arsinh}(g\sqrt{a_0a_1}),\qquad
 c_b=\cosh\gamma\sqrt{a_b/a_{1-b}},\qquad
 \kappa_b=(c_b-1)/g.
\]
Write \(q_n=(\psi_0^{(n)}(0),\psi_1^{(n)}(0))^t\) for the next graph
jets and \(s_n=(S_0^{(n)}(0),S_1^{(n)}(0))^t\).  For \(n\ge3\),
\begin{equation}\label{eq:v25-intro-jet-block}
 s_n=M_nq_n+R_n,\qquad
 M_n=\begin{pmatrix}
 \coth(n\gamma)&\mathfrak r_0^n\operatorname{csch}(n\gamma)\\
 \mathfrak r_1^n\operatorname{csch}(n\gamma)&\coth(n\gamma)
 \end{pmatrix},\qquad
 \mathfrak r_b=\sqrt{c_b/c_{1-b}},\quad\det M_n=1,
\end{equation}
where \(R_n\) is determined by lower graph jets.  Establishing this
filtration requires more than the density cancellation.  We invert the
nonlinear half-line stationarity operator on a weighted sequence space,
differentiate finite truncations of its action, control the endpoint tails,
and isolate the first homogeneous occurrence of each new graph jet.
A finite Taylor-remainder estimate proves that equal graph jets through
order \(M\) give equal action jets through order \(M\), even for
smooth graphs differing by flat terms.  Smooth boundary images are not
identified by their jets; the boundary-image conclusion uses analyticity.
The inverse is fixed-order stable; determinant one alone is not a bound
uniform in the infinite jet order.

For a periodic table the deck group is a marked abstract \(\mathbb Z^2\),
with no supplied Euclidean Gram form.  A channel label \((a,b,\ell)\)
specifies its obstacle orbits and abstract deck displacement.  Each channel
has its own oriented frame; relative placements between channels are not
supplied.  A signature-rigid transition has a unique framed point with the
prescribed complete oriented curvature signature.  In the presence of
analytic symmetries the general gluing theorem gives a gluing space, not
an arbitrarily selected phase.

\begin{theorem}[Determination from single-offset transverse laws]
\label{thm:v26-intro-rigidity}
\label{thm:v25-intro-rigidity}
For a smooth strictly convex channel, its gap and two same-type signed
limiting endpoint laws at one known sufficiently small positive offset
 determine every finite labelled contact jet.  No longitudinal position,
residual time, waiting count, flux amplitude, or contact curvature is
supplied as part of these conditional-law data.

For connected real-analytic incident boundaries the data determine their
complete boundary images in the separate channel frames.  Suppose the
marked data are realizable by a periodic table and the measured channel
graph contains a rank-two anchoring pair based at one channel frame,
in the sense of Definition~\ref{def:v24-rank-two-anchoring}, and a rooted
signature-rigid spanning tree reaching every obstacle orbit.  Then the
data determine the entire marked periodic table,
including its Euclidean lattice, up to one simultaneous element of
\(\mathrm{SE}(2)\).  On a compact class with uniform collars the same
positive offset can be used for every channel and both types.
\end{theorem}

\begin{proof}
Theorem~\ref{thm:v26-density-inverse} supplies the action germs.  The
all-order contact inverse is Theorem~\ref{thm:v22-signed-rigidity}.
Analytic continuation and the intrinsic gluing classification are proved
in Theorem~\ref{thm:v23-intrinsic-table-rigidity}; rank-two metric recovery
is Theorem~\ref{thm:v24-lattice-gram-recovery}.  Propagation along the
rooted spanning tree is Theorem~\ref{thm:v24-uncalibrated-periodic-rigidity}.
To use the same frame for these last two steps, apply
Lemma~\ref{lem:v40-signature-rigid-rerooting} to root the supplied tree
at an obstacle met by the anchoring pair.  A signature-rigid incidence
on that pair supplies the unique framed signature required by the lemma.
Thus the two cycle holonomies and the tree placements are reconstructed
in one common frame, without an additional registration hypothesis.
Their composition is Theorem~\ref{thm:v26-single-offset-global}.
\end{proof}

The last lattice step is \(L=VM^{-1}\), where \(V\) consists of recovered
translation holonomies and \(M\) of their independent marked deck
vectors.  The difficult input is the recovery and unique matching of
complete framed curve images, not this matrix identity.  Likewise, marked
deck labels and signed transverse conventions are part of the intrinsic
data map, not quantities inferred from an unlabelled trajectory.

\subsection{Physical acquisition and the same-experiment comparison}

For global physical acquisition each channel detector reports the first and
last planar collision positions in a common oriented sensor frame for that
channel's two types, together with residual preparation time or a failure
symbol.  Frames of distinct channels are not registered.  Designs use only
the marked channel, type, even flight number, and physical observation time.
All preparations and stopping decisions are retained and charged.  These
common Borel spaces are defined in \eqref{eq:v25-design-record-spaces}.
The anchored scalar local experiment is a different, explicitly coarser
model; no unknown contact chart is silently treated as observed.

Let \(\mathfrak K\) be the compact marked analytic class of
Section~\ref{sec:v25-signature-stability}, with uniform positive geometry,
holomorphic curvature collars, persistent channels, the preceding uniqueness
structure, and nondegenerate rank-two holonomies.  Let \(d_q\) be the
labelled \(C^q\) distance in a root-frame gauge.

\begin{theorem}[Uniform physical reconstruction at one post-pilot offset]
\label{thm:v26-intro-physical}
\label{thm:v25-intro-physical}
There are finite deterministically capped policies \(\pi_s\) and measurable
table estimators \(\widehat T_s\) in the common physical position experiment
such that every flight number used by \(\pi_s\) is even and at least
\(s\), and for every fixed \(q\ge0\) and \(\epsilon>0\),
\[
 \sup_{T\in\mathfrak K}\Pr_T^{\pi_s}
       \{d_q(\widehat T_s,T)>\epsilon\}\longrightarrow0.
\]
No gap, contact center, tangent, inter-channel placement, or lattice metric
is supplied.  After the pilot all channels and both types use one fixed
positive offset from the estimated onset, independent of \(s\).  The
estimator uses the estimated gaps and bounded tests of the successful
transverse endpoint pairs; it does not require their normal coordinates
or residual times in its post-pilot criterion.
\end{theorem}

\begin{proof}
Theorem~\ref{thm:v25-observable-calibration} gives a charged near-onset
pilot at the already selected final even flight number.  The
single-offset inverse makes gaps and the two endpoint laws at a fixed
\(d_0>0\) jointly separating.  Lemma~\ref{lem:v26-single-offset-separators}
selects finitely many bounded Lipschitz tests, and
Theorem~\ref{thm:v26-fixed-order-physical} gives a finite-template estimator
with a deterministic preparation cap.  The compact inverse modulus and
Theorem~\ref{thm:v26-global-physical-reconstruction} give the conclusion.
\end{proof}

The theorem does not claim that half-line inversion is necessary in this
position experiment.  Proposition~\ref{prop:v26-position-jets} gives a
second route using exactly the same long even designs, charged pilot, and
position records.  Once the frame error is at most \(\varepsilon\),
interpolation of \(M+1\) separated graph samples at scale \(a\) gives
\[
 |\widehat\psi_b^{(m)}(0)-\psi_b^{(m)}(0)|
 \le C_M\{a^{M+1-m}+\varepsilon a^{-m}\},\qquad m\le M.
\]
Each required bin has positive preparation probability and a finite charged
cap.  Choosing \(\varepsilon\le a^{M+1}\) gives fixed-order consistency
without estimating an action.  Arbitrarily large even flight numbers do
not exclude this comparison.

The intrinsic theorem has different data: no graph ordinate is available.
Its single-offset distributional inverse and its all-order alternating
contact recursion give a reconstruction from genuinely coarsened data.
Exact law injectivity is not statistical equivalence to noiseless position
sampling.  Proposition~\ref{prop:v26-position-deficiency} identifies
the precise dependence on the observed contact type in a registered
homothetic family.  For any finite number of successful endpoint pairs,
the scalar-to-position deficiency is one when the even design starts
on the varying obstacle, while the reverse deficiency is zero.  When
both observed endpoints instead lie on the fixed facing obstacle, the
matched endpoint experiments are equivalent.  Residual time is deleted
from both records in this comparison; retaining it in both gives the
same dichotomy.  Intermediate collisions are not silently added to either
record.  These comparisons do not alter the two-type law inverse above.

The finite separator library and compact inverse moduli may depend
non-effectively on \(\mathfrak K\).  Neither global estimator asserts an
analytic minimax rate or polynomial-time design construction.  Separately,
Theorem~\ref{thm:v25-finite-signature-embedding} supplies a uniform finite
signature with local embedding and global separation bounds, and
Lemma~\ref{lem:v25-noisy-signature-match} gives unique matching under
controlled \(C^2\) perturbations.  The compact inverse modulus does not
replace that matching proof.  Budget-indexed implementations run the selected
stage afresh; a cumulative archive of earlier stages is not claimed to
have a diverging minimum flight number.

\subsection{Local information under coarsening}

In the anchored finite-dimensional family, residual-time retention leaves
positive density at a moving ceiling.  At shape scale \(k_n^{-1}\), the
successful endpoint--time experiment is equivalent to a Poisson boundary-shift
experiment; the corresponding gap scale is \((j_nk_n)^{-1}\).
Theorem~\ref{thm:v24-two-sided-deficiency} gives parameter-independent kernels
in both directions, controlling the layer, reconstructed bulk, and the
intersection with the fixed residual-time face.  Its bulk estimate uses the
explicit bounded \(O(k_n^{-1})\) relative-density perturbation, not bare
trace convergence.

The finite-strip Poisson family is jointly dominated by a common Poisson
envelope.  It need not be absolutely continuous with respect to its
zero-shift member.  The scalar endpoint family is likewise jointly dominated
by Lebesgue measure with the required cemetery atoms.  Its common-collar
censoring produces the stronger property of domination by the censored
reference member.  Proposition~\ref{prop:v26-common-domination} gives both
reference measures and the Poisson Radon--Nikodym formula explicitly.

Removing residual time produces linear vanishing at the endpoint support
boundary and the scale \(k_n\delta_n^2\log(1/\delta_n)\to1\).
The reference-dominated representative has a uniform LAN expansion.
Lemma~\ref{lem:v33-finite-likelihood} identifies its finite Gaussian limits;
the uniform moving-boundary modulus and finite-net argument then prove
convergence of the original endpoint family in two-sided Le Cam distance
on every fixed compact local set, in Theorem~\ref{thm:v32-compact-vector}
and Corollary~\ref{cor:v32-compact-fixed-window}.
Waiting counts identify the normal quotient to a level set of the
hyperbolic exponent at scale \((j_n\sqrt{k_n})^{-1}\), while endpoint
information acts on the tangent directions at their slower scale.
Theorem~\ref{thm:v23-count-endpoint-joint} gives the compact two-sided
Gaussian product approximation for this count--endpoint record.
Deleting those counts is a further coarsening.  These assertions concern
their declared scalar record spaces, not the support-singular planar
position experiment or the complete growing collision array.

\subsection{Comparison with related inverse problems}

Marked lengths, dynamical conjugacy, and boundary-law determination
pose different inverse problems.  De Simoi, Kaloshin and
Leguil~\cite{DSKL} determine analytic open dispersing billiards from
marked lengths under their symmetry and genericity hypotheses.
Florio and Leguil~\cite{FlorioLeguil} obtain smooth conjugacies from
periodic length data for the hyperbolic systems covered by their
hypotheses, including the stated open-billiard application.  This is a
dynamical conclusion, distinct from identifying obstacle images by a
Euclidean isometry.\footnote{We use the revised version v5 of
\cite{FlorioLeguil}.  Its revision notice removes the earlier
spectral-rigidity assertion affected by the error in Proposition 3.1;
the smooth-conjugacy result is the comparison used here.}
Finamore and Leguil~\cite{FinamoreLeguil} use an enriched marked length
spectrum and a CAT(0) construction to recover finite-horizon Sinai
billiards up to isometry.

The data in Theorem~\ref{thm:v26-intro-rigidity} are instead onsets and
signed transverse laws of selected channels, with marked deck labels
and signature-rigid incidences.  The local observation deletes the
normal endpoint coordinates while retaining the two labelled contact
types.  The proof first identifies its nonlinear limiting density from
physical bridges, and then recovers the unsymmetrized actions and the
actual contact jets.  Only the subsequent boundary-image step uses
analyticity.  This composition does not pass through a supplied
normal-form chart or a dynamical conjugacy.  Neither the channel laws
and incidence data nor the ordinary or enriched marked length spectrum
are derived from the other observation map here; the respective
hypotheses and conclusions therefore remain distinct.

The cofactor identity and inverse-decay estimates used in the forward
argument are classical tools, as detailed in
Section~\ref{sec:v9-operators}.  Their role here is to control a
nonlinear Hessian perturbation in trace norm along the whole physical
bridge, before passing through a normalized moving-sublevel integral.
The inverse then requires finite-truncation envelope differentiation
and a smooth finite-remainder theorem before homogeneous degree
isolation yields the contact recursion.  These are the analytic and
geometric steps connecting the relative law to the determination
theorem; the determinant-one identity and the final lattice equation
record their explicit inversion, rather than replacing those steps.

Within the present geometry, the mixed-type one-flight support inverse of
Proposition~\ref{prop:v25-one-flight-inverse} and the two-flight calculations
remain useful benchmarks.  The direct-position result above is stronger as
a comparison to our physical protocol because it uses the same long even
record restriction.  The additional inverse content lies in recovering
geometry from the single-offset transverse laws, not in prohibiting one
short flight while supplying noiseless geometric samples.

Nonregular support-dependent estimation and exceptional normal rates have
an established history; see Smith~\cite{Smith1985}.  Meister and
Reiss~\cite{MeisterReiss} prove equivalence of a nonregular regression
experiment to Poisson experiments whose intensity supports encode the
regression boundary.  A Gaussian/Poisson contrast or Poissonization alone
is therefore not asserted as a new principle.  The results here identify
the billiard-specific nonlinear density, its amplitude-free and all-order
contact inverses, and the kernels linking explicitly retained observation
levels, with the rare-event preparation cost included.

\subsection{Organization}

Part I proves the relative law, the all-order signed contact inverse, the
single-offset density inversion, intrinsic analytic gluing, lattice recovery,
and finite-signature stability.  It retains the short-flight comparisons
and includes the same-experiment direct-position benchmark.  Part II proves the
boundary information results, distinguishes common domination from reference
absolute continuity, proves the observed-type information dichotomy,
and compares the residual-time and count coarsenings.
Part III constructs observable calibration and the single-offset global
estimator, and separately gives compatible analytic variation bundles.
The appendices retain the complete auxiliary derivations and applications.
